\section{The construction side}\label{sec:construction}

Our aim at the outset was a linear-time transposer. We did not find one. This section records what the
results above require of any construction, describes the one architecture we know that passes every
test, and summarizes the mechanisms that fail. It contains no new theorems: every claim is either one of
the results above, a citation, or explicitly labelled as a heuristic.

\subsection{Necessary conditions}

A machine witnessing~(T) must have all of the following properties.
\begin{enumerate}[label=(N\arabic*),leftmargin=3em]
\item \emph{Coding.} It cannot be a pure data mover with an input-independent pattern: moving the entries as indivisible items along a fixed pattern needs time $\Omega(L\log d)$ (Sto\ss~\cite{Stoss1973} for Turing machines, Aggarwal and Vitter~\cite{AggarwalVitter1988} in
external memory; both bounds concern a fixed permutation of distinct items). For data movement that depends on the entries, which repeat because $b=\clog d$, no such bound is known. Oblivious machines that compute on the data are not covered by these bounds either; whether they can transpose in linear time is open (see (N4) and \cref{thm:oblivious}).
\item \emph{Many passes, at every scale.} Its run has $\Omega(\log d)$ passes (Theorem~P$'$,
\cref{thm:Pprime}), even on a $2^{-O(db)}$ fraction of the inputs (\cref{thm:Pprime-on}); the grouped form of
\cref{thm:Pprime} gives the same at coarser column scales.
\item \emph{Folding motion.} Along the popular trace, the macroscopic work-head trajectories satisfy $(\star)$
(\cref{cor:mstar-T}): some head makes a macroscopic U-turn, and the heads together waste at least one unit of range (\cref{thm:excess}). Moreover, the limiting paths are locally affine only off sets of times whose measures, summed over the heads, are at least $\tfrac12$~\cite{BrendeAstra2026}; schedules whose macroscopic motion is piecewise affine are excluded.
\item \emph{A large coding gain.} The block network of the run of a slightly modified machine (output buffered at
the end; state and head positions carried on extra stationary tapes) carries a causal code for complete exchange,
on a set of co-entropy $O(d)$, at the profile capacities (\cref{thm:reduction,thm:coent-linear}); by
\cref{thm:code-gap} this code beats routing by a factor $\Omega(\log d)$. For an oblivious machine the code works
on all inputs (\cref{thm:oblivious}); for oblivious I/O machines the analogous statement, at the level of pages, is
\cite[Thm.~6.2]{AHJKL2006}.
\item \emph{Not an XOR-frame recursion.} Within the abstract model of \cref{sec:frames}, every run realizing a gadget makes at least as many frame changes as the gadget has wires (\cref{thm:frame-change}), so every gadget has cost $s\ge W$ (\cref{prop:gadget-runs}, not formalized), while linear time would need $s<W$; the
model captures the preprint's template only informally.
\end{enumerate}

Condition~(N4) is the heart of the matter. Coding gains over routing in undirected networks do exist: the
multiple-unicast (network coding) conjecture is false~\cite{BravermanHe2026,ZhangLiLi2026}. But the known growing
gains are of order $(\log n)^{\varepsilon}$ on networks of size $n$, obtained by amplifying a constant gap by tensoring~\cite{BGS2017}, and
they concern sparse demands; we do not expect tensoring a constant gap to produce a gain of order $\log n$ for dense
demands (a heuristic, not a theorem). Braverman and He also observe~\cite{BravermanHe2026} that the unrefereed sub-$n\log n$ Fourier circuits
of~\cite{OpenAI2026fourier} would give a coding gain over routing; such gains would not come from tensoring, but they
live on circuit graphs, not on networks shaped by tapes. A linear-time transposer needs a gain of order $\log d$ for the densest demand, on a network shaped by tapes, which by the heuristic above would have to be \emph{native}, that is, not obtained by amplifying a constant gap.

\subsection{An architecture that survives}

We describe, as a design and not as a result, the one architecture that we could not rule out.

\paragraph{Geometric recursion.} The transposer runs $\Theta(\log d)$ levels. Level $\ell$ works on data of
volume $\theta^\ell L$ for a constant $\theta<1$, with $O(1)$ passes of an encoder $E_\ell$ and a decoder
$D_\ell$: it computes its output as $D_\ell(X_\ell,f_{\ell+1}(E_\ell(X_\ell)))$, where $f_{\ell+1}$ is the next
level's problem. The total time is $\sum_\ell O(\theta^\ell L)=O(L)$, and the pass count is $\Theta(\log d)$, as
Theorem~P requires.

\paragraph{Expanding choreography.} The heads of each level are driven along prepared \emph{rulers} that realize
a constant number of monotone maps of the interval, so that the relation graph between the encoder's input and
output rows is a bounded-degree expander. The monotone expanders of Bourgain and
Yehudayoff~\cite{BourgainYehudayoff2013} provide such maps (near-identity M\"obius maps from a suitable free
generating set, together with translations and the identity); at least three monotone maps are needed, since a union of two monotone matchings is planar~\cite{DSW2016}, and graphs of treewidth $o(d)$, planar ones included, admit no constant-fraction saving~\cite{BrendeAstra2026}. The rulers can be written in linear time, since their marks are
computed with $O(\log L)$-bit arithmetic (a design computation, not included here). Near-identity generators are needed: integer or dyadic M\"obius pairs generate
discrete groups whose relation graphs are not expanders (a heuristic supported by numerical spectral gaps).
Repeating each choreography a constant number of times should make the macroscopic motion satisfy the quantitative inequality $(\star)$; this is a cost constraint, not an obstruction (heuristic). The follow-up~\cite{BrendeAstra2026} proves a data-free version: head motion along rational near-identity M\"obius maps, generated by a fixed machine in linear time, that satisfies the cut inequality and the conclusions of Theorem~M*; coupling such motion to the data remains open. It also excludes contracting recursions whose stages all have regular (piecewise affine) limiting motion, unless irregular stages take a constant fraction of the time. Stages driven along non-affine M\"obius maps are not regular, so the design above escapes; its affine variants do not.

\paragraph{Two coding cores.} The geometry is the easy part. Each level needs a constant-fraction coding saving at
the encoder, so that $\abs{E_\ell(X_\ell)}\le\theta\abs{X_\ell}$, and, because the decoder writes the level's whole
output from $O(1)$-pass views of the next level's output, also at the decoder. In the simplest linear form the
encoder core is the following problem.

\begin{quote}
\emph{Locally encoded index coding with own-row side information.} Let $\Gamma$ be a bounded-degree graph on
$[d]$. A broadcast consists of linear functionals of a $d\times d$ matrix, each supported on the two rows of an
edge of $\Gamma$. Receiver $j$ sees the broadcast and knows row $j$, and must recover column $j$. Can the broadcast
have size $(1-\varepsilon)d^2$ for a fixed $\varepsilon>0$?
\end{quote}

Without the degree bound, $\varepsilon=\tfrac12$ is achievable: broadcast $A+A^{\T}$ above the diagonal. With the degree bound, the follow-up~\cite{BrendeAstra2026} settles two cases negatively. If $\Gamma$ has treewidth $o(d)$ (for instance planar or hyperfinite graphs), no fixed $\varepsilon$ is possible, even for nonlinear and adaptive broadcasts with a helper. On every bounded-degree $\Gamma$, expanders included, broadcasts that use each row only through the coordinates of one recoding of that row, which include routing-type packets, save only a fraction $O(\log\Delta/\log d)$ of the bits. The question remains open for expanders with general two-row functionals. A positive answer on expanders would be, to our knowledge, the first
constant-fraction coded saving in a tape-relevant model. A negative answer, together with a negative answer to the
analogous decoder problem, would rule out every geometric recursion of the form above with $O(1)$-pass levels
(heuristic: the necessity of both savings is argued here only informally).

\subsection{Mechanisms that fail}

The following classes of constructions fail for reasons proved in this paper.
\begin{itemize}[leftmargin=2em]
\item \emph{Few passes.} Any schedule with $o(\log d)$ passes, with any coding, including constant-pass
choreographies along expanders (Theorems~P and~P$'$).
\item \emph{Regular or one-way motion.} Machines whose macroscopic head motions (along the popular trace) are all
stationary or monotone (\cref{thm:excess,thm:mstar-adaptive}).
\item \emph{XOR-frame recursions} in the abstract model of \cref{sec:frames}, which captures the template
of~\cite{OpenAI2026mult} informally: every run realizing a gadget makes at least as many frame changes as the gadget has wires (\cref{thm:frame-change}), so every gadget has cost at least $W$ (\cref{prop:gadget-runs}, not formalized) and linear time is out of reach.
\end{itemize}
Some classes of codes are moreover no easier than plain codes: linear codes, codes correct off a small set, and codes
that are linear outside $O(dw)$ bits of capacity exist only if plain codes exist (\cref{thm:pr-linear,thm:pr-vote});
a counterexample to the robust conjecture from these classes would refute the plain one.
Many further attempts failed for reasons we could only argue informally: flat schemes that code across a constant
number of whole streams (they save at most a constant factor), window-local recursions whose levels read only
nearby data, static coding certificates on Ramanujan-type graphs that resist conversion into causal codes, and
re-indexings that turn transposition into a cheap map only through a transposition-hard change of basis. We mention them as warnings, not as results. The follow-up~\cite{BrendeAstra2026} excludes further classes, with formal proofs: bounded-length words of whole-row operations (integer row addition and subtraction, XOR prefix sums, reversals, rotations, masked strides, and a ternary palette); contracting recursions whose stages all have regular limiting motion, as above; and a two-stage nonlinear architecture, which needs about $L/4$ helper bits.
